\appendixitem{THE CONNECTED COMPONENTS OF A GRAPH}

Let \(\Gamma = (\mathbf{A},\mathbf{S})\) be a graph. If \(a\) and \(b\) are two vertices, a path joining \(a\) and \(b\) is a sequence \((x_0,\ldots ,x_n)\) of vertices of \(\Gamma\) with \(x_0 = a,x_n = b\) , the vertices \(x_{i}\) and \(x_{i + 1}\) being joined for \(0\leqslant i < n\) ; the integer \(n\geqslant 0\) is the length of the path. The path \((x_0,\ldots ,x_n)\) is said to be injective if \(x_{i}\neq x_{j}\) if \(i\neq j\) . If a path \((x_0,\ldots ,x_n)\) joining \(a\) and \(b\) is of minimal length among such paths, it is necessarily injective; for if not, there would exist i and j with \(0 \leqslant i < j \leqslant n\) and \(x_{i} = x_{j}\) and then the sequence

\[
\left(x _ {0}, \dots , x _ {i}, x _ {j + 1}, \dots , x _ {n}\right)
\]

would be a path of length \textless{} n joining a and b.

The relation ``there exists a path joining a and b'' between two vertices a and b of \(\Gamma\) is an equivalence relation R on the set S of vertices. The equivalence classes of R are called the connected components of \(\Gamma\) ; and \(\Gamma\) is said to be connected if S has at most one connected component, that is if any two vertices of \(\Gamma\) can be joined by at least one path.

PROPOSITION 1. Let \(\Gamma = (A, S)\) be a graph and \((S_{\alpha})_{\alpha \in L}\) its family of connected components. Denote by \(\Gamma_{\alpha}\) the full subgraph of \(\Gamma\) having \(S_{\alpha}\) as its set of vertices.

\begin{enumerate}
\def\labelenumi{(\roman{enumi})}
\item
  For any \(\alpha\) in L, the graph \(\Gamma_{\alpha}\) is connected.
\item
  If \(\Gamma' = (A', S')\) is a connected subgraph of \(\Gamma\) , there exists \(\alpha\) in L such that \(\Gamma' \subset \Gamma_{\alpha}\) .
\item
  If \(\alpha \neq \beta\) , no element of \(S_{\alpha}\) is joined in \(\Gamma\) to any element of \(S_{\beta}\) (equivalently, every edge of \(\Gamma\) is an edge of one of the \(\Gamma_{\alpha}\) ).
\item
  Let \((\mathrm{S}_{\lambda}^{\prime})_{\lambda \in \mathbb{M}}\) be a partition of \(S\) such that, if \(\lambda \neq \mu\) , no element of \(\mathrm{S}_{\lambda}^{\prime}\) is joined in \(\Gamma\) to any element of \(\mathrm{S}_{\mu}^{\prime}\) ; then each of the sets \(\mathrm{S}_{\lambda}^{\prime}\) is a union of connected components of \(\Gamma\) .
\item
  Let \(\alpha\) be in L and \(a\) and \(b\) be in \(S_{\alpha}\) . Then there exists a path \(c = (x_0, \ldots, x_n)\) joining \(a\) and \(b\) in \(\Gamma\) . For any \(i\) with \(0 \leqslant i \leqslant n\) , the path \((x_0, \ldots, x_i)\) joins \(a\) to \(x_i\) in \(\Gamma\) , so \(x_i \in S_{\alpha}\) . Thus, \(c\) is a path in \(\Gamma_{\alpha}\) joining \(a\) and \(b\) . It follows that \(\Gamma_{\alpha}\) is connected.
\item
  Let \(\Gamma' = (A', S')\) be a non-empty connected subgraph of \(\Gamma\) , let \(a\) be an element of \(S'\) and let \(S_{\alpha}\) be the connected component of \(S\) containing \(a\) . For any \(b\) in \(S'\) , there exists a path \(c\) joining \(a\) and \(b\) in \(\Gamma'\) , and \(a\) fortiori in \(\Gamma\) . It follows that \(S' \subset S_{\alpha}\) .
\item
  Given distinct elements \(\alpha\) and \(\beta\) of \(L\) , and vertices \(a \in S_{\alpha}\) and \(b \in S_{\beta}\) , there is no path joining \(a\) and \(b\) , and in particular no edge joining \(a\) and \(b\) .
\item
  Let \(a\) be in \(S_{\lambda}'\) and let \(S_{\alpha}\) be the connected component of \(\Gamma\) containing \(a\) . Then, for any \(b\) in \(S_{\alpha}\) , there exists a path \((x_0, \ldots, x_n)\) joining \(a\) and \(b\) in \(\Gamma\) . If \(i\) is an integer such that \(0 \leqslant i < n\) and if \(x_i \in S_{\lambda}'\) , then \(x_{i+1} \in S_{\lambda}'\) since \(x_i\) is joined to \(x_{i+1}\) . It follows by induction that \(x_i \in S_{\lambda}'\) for \(0 \leqslant i \leqslant n\) , and in particular that \(b = x_n\) is in \(S_{\lambda}'\) . Thus, \(S_{\alpha} \subset S_{\lambda}'\) .
\end{enumerate}

COROLLARY 1. A graph \(\Gamma = (A, S)\) is connected if and only if there does not exist a partition \((S', S'')\) of \(S\) into two non-empty subsets such that no element of \(S'\) is joined in \(\Gamma\) to any element of \(S''\) .

Suppose that \(\Gamma\) is not connected and let \(S'\) be one of its connected components. The set \(S'' = S - S'\) is non-empty by Prop. 1, (i) and no element of \(S'\) is joined to any element of \(S''\) by Prop. 1, (iii).

Suppose now that \(\Gamma\) is connected and let \((S', S'')\) be a partition with the stated property. By Prop. 1, (iv), the set \(S'\) contains at least one connected component, so \(S' = S\) and \(S'' = \varnothing\) , a contradiction.

COROLLARY 2. A subset \(S'\) of \(S\) is a union of connected components if and only if no vertex belonging to \(S'\) is joined to any vertex belonging to \(S - S'\) .

The condition is sufficient by Prop. 1, (iv). It is necessary by Prop. 1, (iii).

% semantic-fragments: legacy-input-seam
\relax{}
